\chapter{عزوم دالة زيتا والقيم الكبيرة والمتطرفة}
\label{ch:moments-extreme-values}

\section*{نطاق الفصل وحالته}

يدرس هذا الفصل المتوسطات التكاملية لقيم دالة زيتا على الخط الحرج، ثم يربطها بالحدود العليا والدنيا وبالنموذج الحدسي للمصفوفات العشوائية. النواة الصارمة هي دالة زيتا لريمان؛ أما عائلات دوال \(L\) فتظهر بوصفها امتدادًا مفاهيميًا مضبوطًا لا نقلًا آليًا للثوابت.

\noindent\textbf{حالة الفصل:}
\textenglish{AUTHORED-DRAFT / NON-CITABLE}. اجتازت الحزمة العلمية مراجعة مستقلة ومراجعة ضيقة بحكم \textenglish{PASS-FOR-AUTHORING = YES}. تبقى النتائج غير قابلة للاستشهاد حتى البناء والتدقيق والمراجعة المستقلة اللاحقة واعتماد المالك.

\subsection*{حراس النطاق}
\begin{itemize}
  \item كل حد من Soundararajan أو Harper يذكر فرضية ريمان صراحة.
  \item نتيجة Harper توصف بأنها \textenglish{preprint}، ولا تسوى مكانتها الببليوغرافية بمقالة Soundararajan المحكمة.
  \item العزم الرابع نتيجة مقتبسة من Ingham، لا برهان داخلي كامل.
  \item حدسية Keating--Snaith حدسية، لا مبرهنة.
  \item لا يستنتج قانون للقيمة القصوى النموذجية من عزم منفرد.
\end{itemize}

\subsection*{نواتج التعلم}
بنهاية الفصل يستطيع القارئ أن:
\begin{enumerate}
  \item يميز بين رتبة العزم \(2k\) والرمز \(I_k(T)\).
  \item يستعمل المعادلة الوظيفية التقريبية وصيغة القيمة المتوسطة لكثيرات حدود ديريشليه.
  \item يفهم منشأ الحدين الرئيسيين في العزم الثاني.
  \item يقرأ صيغة Ingham للعزم الرابع مع تصنيفها الصحيح.
  \item يفرق بين الحدود العليا المشروطة والحدود الدنيا غير المشروطة.
  \item يفصل العامل الحسابي عن عامل المصفوفات العشوائية في حدسية العزوم.
  \item يحدد بدقة ما تعطيه متراجحة ماركوف عن الذيول وما لا تعطيه عن القيم القصوى.
\end{enumerate}

\section{تعريف العزوم وتطبيعها}

\begin{definition}[العزم المستمر على الخط الحرج]
\label{def:ch22-continuous-moment}
\resultid{ANT-DEF-22-01}
\provedhere
لكل \(k\ge0\) نضع
\[
I_k(T)=\int_0^T\left|\zeta\!\left(\frac12+it\right)\right|^{2k}\,dt,
\qquad
J_k(T)=\int_T^{2T}\left|\zeta\!\left(\frac12+it\right)\right|^{2k}\,dt.
\]
الرتبة هي \(2k\)، لا \(k\). ويستعمل \(J_k(T)\) في النتائج المحلية على نافذة ديادية، بينما يستعمل \(I_k(T)\) في الصيغ التراكمية.
\end{definition}

العزوم تقيس توزيع الكتلة لا القيمة القصوى وحدها. فإذا كان \(V>0\)، فإن متراجحة ماركوف تعطي
\[
\operatorname{meas}\left\{t\in[T,2T]:|\zeta(1/2+it)|\ge V\right\}
\le V^{-2k}J_k(T),
\]
لكنها لا تحدد شكل الذيل الدقيق ولا موضع القمة النموذجية.

\section{المعادلة الوظيفية التقريبية}

\begin{proposition}[صيغة تقريبية متناظرة لدالة زيتا]
\label{prop:ch22-afe}
\resultid{ANT-PROP-22-01}
\provedhere
إذا كان \(t\ge2\)، واختير \(x,y>0\) بحيث \(xy=t/(2\pi)\)، فإن
\[
\zeta\!\left(\frac12+it\right)
=
\sum_{n\le x}\frac1{n^{1/2+it}}
+
\chi\!\left(\frac12+it\right)
\sum_{n\le y}\frac1{n^{1/2-it}}
+O\!\left(x^{-1/2}+y^{-1/2}\right),
\]
حيث
\[
\chi(s)=\pi^{s-1/2}\frac{\Gamma((1-s)/2)}{\Gamma(s/2)},
\qquad |\chi(1/2+it)|=1.
\]
وبالاختيار المتناظر \(x=y=\sqrt{t/(2\pi)}\) نحصل على مجموعين بطول يقارب \(\sqrt t\)، وحد خطأ من رتبة \(O(t^{-1/4})\).
\end{proposition}

\begin{proof}
نبدأ من المعادلة الوظيفية \(\zeta(s)=\chi(s)\zeta(1-s)\)، ثم نطبق صيغة بيرون المبتورة على السلسلتين في نصفي المستوى المناسبين، ونحرك المسار مع التقاط القطب عند \(s=1\). يوازن الشرط \(xy=t/(2\pi)\) طولي المجموعين، ويعطي تقدير ستيرلنغ لعامل غاما الحدين الباقيين \(O(x^{-1/2})\) و\(O(y^{-1/2})\). هذه هي الصيغة الكلاسيكية المطلوبة فقط؛ لا نحتاج هنا إلى الأوزان الناعمة أو أفضل بقايا معروفة.
\end{proof}

\section{قيمة متوسطة لكثيرات حدود ديريشليه}

\begin{proposition}[صيغة القيمة المتوسطة]
\label{prop:ch22-dirichlet-polynomial-mean}
\resultid{ANT-PROP-22-02}
\provedhere
لتكن \(a_n\in\mathbb C\)، ولتكن
\[
A(t)=\sum_{n\le N}a_n n^{-it}.
\]
عندئذ
\[
\int_T^{2T}|A(t)|^2\,dt
=
T\sum_{n\le N}|a_n|^2
+O\!\left(\sum_{\substack{m,n\le N\\m\ne n}}
\frac{|a_m a_n|}{|\log(m/n)|}\right).
\]
وبالصيغة القياسية لمونتغمري--فوغان:
\[
\int_T^{2T}|A(t)|^2\,dt
=
\left(T+O(N)\right)\sum_{n\le N}|a_n|^2
\]
ضمن النسخة المناسبة للنواة الحالية.
\end{proposition}

\begin{proof}
بالتوسيع نحصل على
\[
|A(t)|^2=\sum_{m,n\le N}a_m\overline{a_n}e^{-it\log(m/n)}.
\]
يعطي القطر \(m=n\) الحد \(T\sum|a_n|^2\). أما خارج القطر فنكامل الدالة الأسية مباشرة، أو نستعمل الجمع الجزئي مع لمّة كوسمين--لانداو واختبار المشتقة الأولى من الفصل الثامن عشر. يعالج المجال القريب من القطر بتقسيم ديادي في \(|m-n|\)، فلا نفترض إلغاء نقطيًا غير موجود.
\end{proof}

\section{العزم الثاني}

\begin{theorem}[العزم الثاني لدالة زيتا]
\label{thm:ch22-second-moment}
\resultid{ANT-THM-22-01}
\provedhere
عندما \(T\to\infty\)،
\[
I_1(T)
=
T\log\frac{T}{2\pi}+(2\gamma-1)T
+O\!\left(T^{1/2}\log T\right).
\]
\end{theorem}

هذه الصيغة الثنائية الحد مع حد الخطأ الكلاسيكي \(O(T^{1/2}\log T)\) تُنسب إلى Ingham؛ انظر \cite{ingham1928meanvalue}. أما أعمال Hardy--Littlewood فتمثل الخلفية التاريخية للحد الرئيس والمعادلة الوظيفية التقريبية، وتعرضها المراجع المعيارية اللاحقة بصيغ حديثة، ومنها \cite{titchmarshHeathBrown1986zeta}.

\begin{proof}
نطبق الصيغة التقريبية المتناظرة بطول \(X(t)=\sqrt{t/(2\pi)}\). لأن \(|\chi(1/2+it)|=1\)، فإن مربعي المجموعين يعطيان حدين قطريين متماثلين، بينما تضبط الحدود المتقاطعة وغير القطرية باستعمال القضية السابقة والجمع الجزئي واختبارات الفصل الثامن عشر.

الحد القطري المركزي هو مجموع توافقي:
\[
\sum_{n\le X}\frac1n=\log X+\gamma+O(X^{-1}).
\]
وبما أن \(X(t)^2=t/(2\pi)\)، فإن جمع نصفي الصيغة يعطي \(\log(t/(2\pi))+2\gamma\). وبعد التكامل نستعمل
\[
\int_1^T\log\frac{t}{2\pi}\,dt
=T\log\frac{T}{2\pi}-T+O(1),
\]
فتظهر
\[
T\log\frac{T}{2\pi}+(2\gamma-1)T.
\]
أما مجموع الحدود غير القطرية وبقايا القطع فيدخل في
\(O(T^{1/2}\log T)\). هذا الحد كلاسيكي وكاف لبنية الفصل، ولا ندعي أنه أفضل حد معروف.
\end{proof}

\section{العزم الرابع: صيغة Ingham}

\begin{theorem}[Ingham، العزم الرابع]
\label{thm:ch22-fourth-moment}
\resultid{ANT-THM-22-02}
\citedresult[\cite{ingham1928meanvalue}]
لدينا
\[
I_2(T)
=
\frac{1}{2\pi^2}T\log^4T+O(T\log^3T).
\]
وبصياغة أدق توجد كثيرة حدود \(P_4\) من الدرجة الرابعة بحيث
\[
I_2(T)=TP_4(\log T)+E_2(T),
\]
ومعامل \(x^4\) في \(P_4(x)\) هو \(1/(2\pi^2)\).
\end{theorem}

لا يعاد بناء برهان Ingham كاملًا هنا. الفكرة البنيوية أن تربيع تقريب زيتا مرتين يولد معاملات دالة القواسم \(d(n)\)، وأن
\[
\sum_{n\le x}\frac{d(n)^2}{n}
\]
ينمو بكثيرة حدود لوغاريتمية من الدرجة الرابعة. لكن استخراج الثابت وضبط الحدود غير القطرية يحتاجان تحليلًا أدق من نطاق هذا الفصل.

\section{حدود عليا عامة تحت فرضية ريمان}

\begin{theorem}[Soundararajan، حد علوي مشروط]
\label{thm:ch22-soundararajan}
\resultid{ANT-THM-22-03}
\citedresult[\cite{soundararajan2009moments}]
نفترض فرضية ريمان. لكل \(k>0\) ثابت ولكل \(\varepsilon>0\)،
\[
J_k(T)\ll_{k,\varepsilon}T(\log T)^{k^2+\varepsilon}.
\]
والحالة \(k=0\) تافهة ويمكن ضمها اصطلاحيًا.
\end{theorem}

يعتمد المسار على حد علوي لـ\(\log|\zeta(1/2+it)|\) بواسطة كثير حدود أولي قصير خارج مجموعة استثنائية مضبوطة، ثم يدمج تقديرات القيم الكبيرة. الشرط RH جزء من المبرهنة وليس ملاحظة جانبية.

\begin{theorem}[Harper، إزالة خسارة \(\varepsilon\)]
\label{thm:ch22-harper}
\resultid{ANT-THM-22-04}
\citedresult[\cite{harper2013sharp}]
نفترض فرضية ريمان. لكل \(k\ge0\) ثابت،
\[
J_k(T)\ll_k T(\log T)^{k^2}.
\]
المصدر \textenglish{preprint} على \textenglish{arXiv:1305.4618}، ولا يعرض هنا بوصفه مقالة محكمة.
\end{theorem}

التحسين حاد من جهة رتبة قوة اللوغاريتم، لكنه لا يثبت التكافؤ التقاربي ولا ثابت Keating--Snaith.

\section{الحدود الدنيا غير المشروطة}

\begin{theorem}[حدود دنيا مستمرة للعزوم]
\label{thm:ch22-lower-bounds}
\resultid{ANT-THM-22-05}
\citedresult[\cite{radziwillSoundararajan2013lower,heapSoundararajan2022lower}]
لكل \(k\ge1\) ثابت، توجد قيمة موجبة \(c_k\) بحيث، لكل \(T\) كبير بما يكفي،
\[
I_k(T)\ge c_k T(\log T)^{k^2}.
\]
النتيجة غير مشروطة، وهي حد أدنى من الرتبة الصحيحة؛ لا تعني
\[
I_k(T)\sim C_kT(\log T)^{k^2}.
\]
\end{theorem}

تظهر هنا قوة طريقة الرنان والمتراجحات الهولدرية: يمكن إثبات الكتلة الدنيا الصحيحة دون امتلاك الصيغة التقاربية الكاملة.

\section{حدسية Keating--Snaith}

\begin{conjecture}[حدسية العزوم العامة]
\label{conj:ch22-keating-snaith}
\resultid{ANT-CONJ-22-01}
لكل \(k\) ثابت مع \(\Re k>-1/2\)،
\[
I_k(T)
\sim
a(k)\frac{G(k+1)^2}{G(2k+1)}
T\left(\log\frac{T}{2\pi}\right)^{k^2},
\]
حيث
\[
a(k)=\prod_p\left(1-\frac1p\right)^{k^2}
\sum_{m=0}^{\infty}
\left(\frac{\Gamma(m+k)}{m!\Gamma(k)}\right)^2p^{-m}.
\]
\end{conjecture}

العامل \(a(k)\) حسابي محلي، بينما
\[
g_{KS}(k)=\frac{G(k+1)^2}{G(2k+1)}
\]
قادِم من متوسطات كثيرات الحدود المميزة في المجموعة الوحدوية. عند \(k=1\)، لدينا \(a(1)=1\) و\(g_{KS}(1)=1\)، فيتوافق الحد مع العزم الثاني. وعند \(k=2\)،
\[
a(2)=\prod_p(1-p^{-2})=\frac6{\pi^2},
\qquad
g_{KS}(2)=\frac1{12},
\]
ومن ثم \(a(2)g_{KS}(2)=1/(2\pi^2)\)، وهو ثابت Ingham. هذا فحص اتساق لا برهان للحدسية.

\section{القيم الكبيرة والمتطرفة}

\begin{openproblem}[القيمة القصوى على النوافذ الطويلة]
\label{open:ch22-extreme-values}
\resultid{ANT-OPEN-22-01}
\openresult
تحديد الرتبة النموذجية الدقيقة لـ
\[
\max_{t\in[T,2T]}\left|\zeta\!\left(\frac12+it\right)\right|
\]
وقانون تقلباتها مسألة مستقلة عن إثبات حد لعزم ثابت منفرد. نماذج الحقول اللوغاريتمية المترابطة والمصفوفات العشوائية تقترح بنية أدق، لكنها لا تتحول تلقائيًا إلى مبرهنة لدالة زيتا.
\end{openproblem}

من حد للعزم نحصل فقط على حد قياس:
\[
\operatorname{meas}\{t:|\zeta(1/2+it)|\ge V\}
\ll V^{-2k}T(\log T)^{k^2+o(1)}.
\]
اختيار \(k\) متغيرًا مع \(V\) قد يعطي معلومات قوية عن الذيول، لكنه يحتاج إلى انتظام في الثوابت ومجال \(k\)، وهي معلومات لا توفرها صيغة ثابتة لـ\(k\) وحدها.

\section{الانتقال إلى عائلات دوال \(L\)}

في عائلة من دوال \(L\)، يجب تحديد ثلاثة أمور قبل كتابة أي حدسية عزوم:
\begin{enumerate}
  \item العائلة نفسها ومقياس عد عناصرها؛
  \item الموصل التحليلي الذي يحل محل \(T\)؛
  \item نوع التناظر: وحدوي أو تعامدي أو سمبلكتي.
\end{enumerate}
لذلك لا ينقل الأس \(k^2\) أو العامل المصفوفي بصورة عمياء بين العائلات. الفصول السابع والحادي والعشرون توفران لغة الموصل والعوامل المحلية، لكن ثوابت العزوم تحتاج تحليلًا خاصًا بكل عائلة.

\section*{خلاصة الفصل}

الفكرة المركزية أن العزوم تقع بين التحليل الصارم والنمذجة الحدسية. العزمان الثاني والرابع مثبتان، والحدود الدنيا العامة غير مشروطة، والحدود العليا الحادة المعروضة هنا مشروطة بفرضية ريمان، أما الصيغة العامة بثابتها الكامل فتبقى حدسية. كذلك فإن الانتقال من العزوم إلى القيم المتطرفة ليس خطوة شكلية، بل جبهة بحثية مستقلة.

\section*{تمارين}
\begin{enumerate}
  \item تحقق مباشرة من \(a(1)=1\).
  \item أثبت أن \(a(2)=1/\zeta(2)=6/\pi^2\).
  \item استنتج من متراجحة ماركوف حدًا لقياس مجموعة \(|\zeta|\ge V\) باستعمال حد Harper.
  \item بين لماذا لا يكفي الحد السابق لتحديد القيمة القصوى النموذجية.
  \item طبق صيغة القيمة المتوسطة على \(a_n=n^{-1/2}\) وفسر ظهور \(\log N\).
\end{enumerate}
